\section{Base loci of divisors}


% Base loci of divisors play an important role in the study of positivity of divisors.
% At least for me, there are two main reasons for this, respectively from \cref{thm:semi-ample_fibration} and \cref{prop:properties_augmented_and_restricted_base_loci}.
% \Yang{}


\subsection{Base loci of line bundles}

  % On this subsection, we do not require $X$ to be projective unless otherwise specified.

  We begin from the support of a Cartier divisor.

  \begin{definition}
    Let $X$ be a variety and $D$ an effective divisor on $X$.
    The support of $D$ is defined as 
    % \[ \Supp D \coloneqq \bigcup_{\xi \in X, f_\xi \notin \calO_{X,\xi}^*} \overline{\{\xi\}}. \]
    \[ \Supp D \coloneqq \{ \xi \in X \mid f_\xi \notin \calO_{X,\xi}^* \}. \]
    Here $\xi$ takes over all scheme-theoretic points of $X$ and $f_\xi$ is a local equation of $D$ at $\xi$, which is well-defined up to multiplication by a unit in $\calO_{X,\xi}$, so the above definition is well-defined.
  \end{definition}

  When $X$ is normal, the support of $D$ is just the union of all prime divisors that appear in $D$ with non-zero coefficients.

  Suppose that $D$ is effective.
  Let $\{U_i\}$ be an open cover of $X$ such that $D$ is defined by $\{U_i, f_i\}$ with $f_i \in \calO_X(U_i)$.
  By taking $U_i$ to be affine, one can see that $\Supp D$ is the union of the closed subsets defined by $f_i$ on $U_i$.
  Hence $\Supp D$ is a closed subset of $X$.
  Furthermore, by Krull's principal ideal theorem, any irreducible component of $\Supp D$ has codimension at most $1$ in $X$.

  In general, let $\xi \in X$ and $f_\xi$ be a local equation of $D$ at $\xi$.
  If $f = f_\xi \in \calO_{X,\xi}^*$, then there exists an open neighborhood $U$ of $\xi$ such that $f \in \calO_X(U)$ defines $D$ on $U$ and is a unit in $\calO_X(U)$.
  This implies that $X \setminus \Supp D$ is open and hence $\Supp D$ is closed.
  
  \begin{remark}
    However, when $X$ is not normal and $D$ is not effective, $\Supp D$ may contain some components of codimension strictly greater than $1$; see \cref{eg:support_of_Cartier_divisor_when_nonnormal_by_gluing_two_points}.
    % \Yang{}
  \end{remark}

  \begin{example}\label{eg:support_of_Cartier_divisor_when_nonnormal_by_gluing_two_points}
    Let $A = \{f \in \kk[s,t] \mid f(0,0) = f(1,0)\}$.
    Geometrically, $A$ is the coordinate ring of variety obtained by gluing two points of $\bbA^2$.
    Let $\frakm_0 = (s,t)$ and $\frakm_1 = (s-1,t)$ be the two maximal ideals of $\kk[s,t]$ corresponding to the two points.
    As $\kk$-modules, we have $\kk[s,t] = \kk \oplus \frakm_0 = \kk \oplus \frakm_1$.
    Then $A = \kk \oplus (\frakm_0 \cap \frakm_1)$ as $\kk$-modules.
    We have $\frakm_0 \cap \frakm_1 = (t, s(s-1))$ (as $\kk[s,t]$-module).
    Set $u = s(s-1)$, $v = s^2(s-1)$ and $w = ts$.
    We have $s^nt = (s^2-s)s^{n-2}t + s^{n-1}t$ and $st,t \in \kk[t,u,v,w]$.
    Hence inductively we have $s^{n}t \in \kk[t,u,v,w]$.
    Similarly, the equation $(s^2-s)s^{n} = (s^2-s)s^{n-2}u+s^{n-1}u$ implies $(s^2-s)s^{n} \in \kk[t,u,v,w]$ since $u,su = v \in \kk[t,u,v,w]$.
    Hence $A = \kk[t,u,v,w] \subset \kk[s,t]$ is a noetherian integral domain but not normal.

    The normalization of $A$ is $\kk[s,t]$ since $s$ is integral over $A$ ($s^2-s-u=0$).
    Moreover, $\kk[s,t,1/u]=A[1/u]$ and $\kk[s,t,1/t]=A[1/t]$ since $s=v/u = w/t$.
    Hence $\Spec \kk[s,t] \setminus V((u,t)) = \Spec A \setminus V((u,t))$.
    Note that $(u,t) = \frakm_0 \cap \frakm_1$ (on both $\kk[s,t]$).
    And on $A$, $A/(u,t) = \kk$ and hence $(u,t) \subset A$ is a maximal ideal.

    Suppose that $\# \kk > 2$.
    Consider the principle Cartier divisor $D = \divisor (s-a)$, $a \neq 0,1 \in \kk$.
    On the open subset $A \setminus V((u,t))$, we can find support of $D$ by looking on $\Spec \kk[s,t] \setminus V((u,t))$.
    Hence $(\Supp D)|_{\Spec \kk[s,t] \setminus V((u,t))} = V(s-a)$.
    On $\frakm = (u,t) \subset A$, note $s \notin A_{\frakm}$ and hence $s-a \notin A_\frakm$.
    Then $\frakm \in \Supp D$.
    Then $\Supp D = V(s-a) \cup \{\frakm\}$.
    We see that $V(s-a) \cong \bbA^1$ is a component of codimension $1$ and $\{\frakm\}$ is an isolated component of codimension $2$.
  \end{example}

  \begin{definition}
    Let $D$ be a (linearly effective) divisor on a variety $X$.
    The \emph{base locus} of $D$ is defined as
    \[ \Bs(D) \coloneqq \bigcap_{D \sim D' \geq 0} \Supp D'. \]
    If $D$ is not linearly equivalent to any effective divisor, we define $\Bs(D) \coloneqq X$.
    Since the definition only depends on the linear equivalence class of $D$, we can also define the base locus for linear equivalence classes of divisors.
  \end{definition}

  This definition gives a closed subset of $X$.
  We can also define the base locus as a subscheme of $X$.

  % \begin{remark}
  %   Due to personal preference, I will try to avoid using the language of linear systems, so I will not use the notation $|D|$.
  % \end{remark}

  \begin{definition}
    Let $D$ be a (linear equivalence class of) divisor on a variety $X$.
    We define the \emph{base ideal} of $D$ as the image of the evaluation map:
    \[ \frakb(D) \coloneqq \Image\Bigl( H^0(X, D) \otimes_\kk \calO_X(-D) \to \calO_X \Bigr). \]
    We define the \emph{scheme-theoretic base locus} of $D$ as the closed subscheme of $X$ defined by the base ideal of $D$.
  \end{definition}

  Let $\pi: X \to \Spec \kk$ be the structure morphism.
  Note that 
  \[ H^0(X, D) \otimes_\kk \calO_X(-D) = \pi^*\pi_*\calO(D) \otimes_{\calO_X} \calO_X(-D), \]
  so the image of the evaluation map is indeed an ideal sheaf on $X$.

  \begin{definition}
    We say a (linear equivalence class of) divisor $D$ on a variety $X$ is \emph{base point free} if $\Bs(D) = \emptyset$.
  \end{definition}

  We also consider the following asymptotic versions of base loci.
  \begin{definition}
    Let $D$ be a (linear equivalence class of) $\bbQ$-divisor on a variety $X$.
    We define the \emph{stable base locus} or \emph{$\bbQ$-base locus} of $D$ as
    \[ \bfB(D) \coloneqq \bigcap_{\substack{m \in \bbZ_{>0},\\ mD \in \Div(X)}} \Bs(mD) = \bigcap_{\substack{D' \in \Div(X)_\bbQ, \\ D \sim_\bbQ D'\geq 0}} \Supp D' \eqqcolon \Bs_\bbQ(D). \]
  \end{definition}

  \begin{definition}
    Let $D$ be a (linear equivalence class of) $\bbQ$-divisor on a variety $X$.
    We say that $D$ is \emph{semi-ample} if $\bfB(D) = \emptyset$.
  \end{definition}

  
  \begin{example}
    Let $X = \Bl_p \bbP^2$ be the blow-up of $\bbP^2$ at a point $p$, $\tilde{H}$ be the strict transform of a line in $\bbP^2$ passing through $p$, and let $E$ be the exceptional divisor.
    Then for $a,b \in \bbZ_{\geq 0}$, we have
    \[ \Bs(a\tilde{H} + bE) = \begin{cases}
      \emptyset & a \geq b, \\
      E & a < b.
    \end{cases} \]
    % \Yang{}
  \end{example}

  \begin{example}[Base locus of divisors on curves]
    Let $C$ be a curve over $\kk$ and $D$ a divisor on $C$.


    Let $X = V(y^2=(x-a_1)\cdots (x-a_n)) \subset \bbP^2$ be a hyperelliptic curve, where $a_i \in \kk$ are point-wise different.
    It is well-known that $X$ is of genus $g = (n+1)(n+2)/2$ and admits a double cover $f$ to $\bbP^1$ via $\kk(x) \subset \kk(x)[y]/(y^2-(x-a_1)\cdots (x-a_n))$.


    base locus on hyperelliptic curve.

    \Yang{}
  \end{example}

  \begin{example}
    Let $E$ be an elliptic curve over $\kkk$ and $P \in E(\kkk)$.
    Set $X = \bbP(\calO_E \oplus \calO_E(-P))$ with the natural projection $\pi: X \to E$ and the unique minimal section $C_0$ with $C_0^2 = -1$.
    Let $D = C_0 + \pi^*P$.
    Then 
    \[ \Bs D = \{Q\}, \]
    where $Q$ is some point on $X$ projecting to $P$ and does not contained in $C_0$.

    Indeed, we have 
    \[ h^0(X, D) = h^0(E, \pi_*\left(\calO_X(C_0) \otimes \pi^*\calO_E(P) \right)) = h^0(E, \calO_E \oplus \calO_E(P)) = 2. \]
    Thus $\dim |D| = 1$.
    For $D_t \in |D|$, write $D_t = H_t + \pi^*B$ where $H_t$ is the horizontal part and $B \in \Div(E)$ is the vertical part.
    Intersecting with fibers, we get that $H_t$ is prime and a section of $\pi$.
    By Bernstein's theorem, for general $t$, $D_t$ is irreducible and thus $D_t = H_t$.
    Since $D^2 = 1$, we only need to find the intersection of two distinct members of $|D|$ to find the base locus.
    Note on $\CH_0(X)$, we have 
    \[ D_t \cap D_s = (C_0 + \pi^*P) \cdot (C_0 + \pi^*P) = [Q] \]
    for a point $Q$.
    By proper pushforward of $\CH_0$, we see that $Q$ projects to $P$.
    It follows that any two irreducible members of $|D|$ intersect at $Q$ since they are section and $\Bs |D| \subset \{Q\}$.
    And $Q \not\in C_0$ since $D_t \cdot C_0 = 0$ for general irreducible $D_t$.
    
    Fix an irreducible member $D_0 \in |D|$.
    For every $D_t = H_t + \pi^*B_t \in |D|$, $H_t$ is different from $D_0$ and thus 
    \[ D_t \cap D_0 = [H_t \cap D_0] + [\pi^*B_t \cap D_0] = [Q]. \]
    Note both $[H_t \cap D_0]$ and $[\pi^*B_t \cap D_0]$ are effective, we get that $[\pi^*B_t \cap D_0] = [Q]$ or $0$.
    If $[\pi^*B_t \cap D_0] = [Q]$, then $B_t = P$ and thus $D_t = H_t + \pi^*P \sim C_0 + \pi^*P$.
    Then $H_t \sim C_0$ and thus $H_t = C_0$.
    It pass through $Q$.
    Otherwise, $[\pi^*B_t \cap D_0] = 0$ and thus $B_t = 0$ and thus $D_t$ is irreducible.
    We have seen that $D_t$ also pass through $Q$.
    Hence, $\Bs |D| = \{Q\}$.

    And we have 
    \[ \Bs(2D) = \emptyset. \]
    This is because $2D = 2C_0 + 2\pi^*P \sim 2C_0 + \pi^*(P_1+P_2)$ for some $P_1,P_2 \in E(\kkk)$ different from $P$.
    Then $Q$ is not contained in $C_0 \cup \pi^{-1}(P_1) \cup \pi^{-1}(P_2)$, so $Q$ is not contained in the base locus of $2D$.
    % \Yang{}
  \end{example}

  We list some basic properties of base loci and base ideals of divisors.

  \begin{proposition}
    Let $D$ be a divisor on a variety $X/\kk$.
    \begin{enumerate}
      \item The underlying closed subset of the scheme-theoretic base locus of $D$ is exactly the base locus of $D$, i.e. 
        \[ \Supp \calO_X/\frakb(D) = \Bs D. \]
      \item For any other divisor $E$ on $X$, we have $\frakb(D) \cdot \frakb(E) \subseteq \frakb(D+E)$ and hence 
      \[ \Bs(D+E) \subseteq \Bs D \cup \Bs E. \]
      \item If $m \mid n$, then $\frakb(mD) \subseteq \frakb(nD)$ and hence $\Bs nD \subseteq \Bs mD$.
      \item $\bfB(D) = \Bs (mD)$ for $m$ sufficiently divisible.
      \item $D$ is base point free on $X \setminus \Bs D$ and semi-ample on $X \setminus \bfB(D)$.
      \item Let $\pi: \tilde{X} \to X$ be a proper birational morphism such that $\pi^{-1} \frakb(D) \cdot \calO_{\tilde{X}}$ is invertible, namely, $\pi^{-1} \frakb(D) \cdot \calO_{\tilde{X}} = \calO_{\tilde{X}}(-F)$ for some effective Cartier divisor $F$ on $\tilde{X}$.
      Then $\pi^* D - F$ is base point free.
      % \item Suppose that $D$ is linearly effective, i.e., $\frakb(D) \neq 0$.
      % Let $\pi: \tilde{X} \to X$ be the blow-up of $X$ along $\frakb(D)$ and $F$ be the exceptional divisor, i.e. $\calO_{\tilde{X}}(-F) = \pi^{-1} \frakb(D) \cdot \calO_{\tilde{X}}$.
      % Then $\pi^* D - F$ is base point free.
      %  \Yang{}
      % \item If $\characteristic \kk = 0$, then there exists a birational morphism $\pi: \tilde{X} \to X$ from a smooth projective variety $\tilde{X}$ such that 
      % \[ \pi^* D = M + F, \quad \Bs M = \emptyset, \quad F \geq 0, \quad \Supp \pi^* D = \Supp F. \]
    \end{enumerate}
  \end{proposition}
  \begin{proof}
    For (a), we check the equality locally.
    Let $\xi$ be a scheme-theoretic point of $X$, $\frakm$ the maximal ideal of $\calO_{X,\xi}$ and $f$ the local equation of $D$ near $\xi$.
    Then $\frakb(D)_\xi$ is generated by $\{gf \mid g \in H^0(X, D)\}$.
    Hence $\frakb(D)_\xi \subseteq \frakm$ if and only if $gf \in \frakm$ for any $g \in H^0(X, D)$, which is equivalent to $\xi \in \Bs D$.

    For (b). 
    The homomorphism $H^0(X, D) \otimes_\kk H^0(X, E) \to H^0(X, D+E)$ given by multiplication of sections induces a homomorphism $\frakb(D) \ten_\kk \frakb(E) \to \frakb(D+E)$ given by the multiplication of sections, so its image is $\frakb(D) \cdot \frakb(E)$ and is contained in $\frakb(D+E)$.
    The inclusion $\Bs(D+E) \subseteq \Bs D \cup \Bs E$ follows immediately from the inclusion of ideals and assertion (a).

    Applying (b) repeatedly we get (c).

    Assertion (d) follows from (c) and noetherianity of $X$.

    For (e), the first assertion follows from the definition of base locus, and the second assertion follows from the first and (d).
    
    For (f).
    By definition, we have a surjective morphism $H^0(X, D) \ten_\kk \calO_X(-D) \surjmap \frakb(D)$.
    Pulling back by $\pi$ and taking the image in $\calO_{\tilde{X}}$, we get a surjective morphism
    \[ H^0(X, D) \ten_\kk \calO_{\tilde{X}}(-\pi^* D) \surjmap \pi^{-1} \frakb(D) \cdot \calO_{\tilde{X}} = \calO_{\tilde{X}}(-F). \]
    That is, 
    \begin{equation}\label{eq:base_ideal_pull_back_by_blow_up}
      H^0(X, D) \ten_\kk \calO_{\tilde{X}} \surjmap \calO_{\tilde{X}}(\pi^* D-F).
    \end{equation}
    % \[ H^0(X, D) \ten_\kk \calO_{\tilde{X}} \surjmap \calO_{\tilde{X}}(\pi^* D-F). \]
    Taking the global sections, we get a map $H^0(X, D) \to H^0(\tilde{X}, \pi^* D-F)$.
    Then pulling back it by $\pi$ and composing with evaluation map, we get
    \begin{equation}\label{eq:base_ideal_pull_back_by_blow_up_and_evaluation}
      H^0(X, D) \ten_\kk \calO_{\tilde{X}} \to H^0(X, \pi^* D-F) \ten_\kk \calO_{\tilde{X}} \to \calO_{\tilde{X}}(\pi^* D-F).
    \end{equation}
    Regarding $\calO_{\tilde{X}}(\pi^* D-F)$ as a subsheaf of the constant sheaf of rational functions on $\tilde{X}$, both the morphism in \cref{eq:base_ideal_pull_back_by_blow_up} and the morphism in \cref{eq:base_ideal_pull_back_by_blow_up_and_evaluation} are given by the same rule, so they are the same morphism.
    Hence, the surjectivity of the morphism in \cref{eq:base_ideal_pull_back_by_blow_up} implies the surjectivity of the evaluation map in \cref{eq:base_ideal_pull_back_by_blow_up_and_evaluation}, so $\pi^* D - F$ is base point free.
    % \Yang{}
  \end{proof}

  \begin{remark}
    In general, we do NOT have $\Bs(mD) \subseteq \Bs (nD)$ for any $m < n$. 
    For example, if $\calO(D)$ is a non-trivial torsion line bundle, then $\Bs(mD) = \emptyset$ for $m$ divisible by the order of $D$ in $\Pic(X)$, and $\Bs(mD) = X$ otherwise.
  \end{remark}

  From now, we consider the asymptotic behavior of divisors on proper varieties.

  \begin{proposition}\label{thm:semi-ample_fibration}
    Let $D$ be a semi-ample $\bbQ$-divisor on a proper variety $X$.
    Then there exists a fibration $f: X \to Y$ to a projective variety $Y$ and an ample $\bbQ$-divisor $A$ on $Y$ such that $D \sim_\bbQ f^* A$.
    Such a fibration is unique up to isomorphism, and it is called the \emph{semi-ample fibration} associated to $D$.
  \end{proposition}
  \begin{proof}
    \hspace{1em}

    \begin{step}
      Existence.
    \end{step}

    Replace $D$ by a sufficiently divisible multiple, we may assume that $D$ is a base point free Cartier divisor.
    % By noetherianity of $X$, there exists finitely many sections $s_0, \dots, s_n \in H^0(X, D)\subset \fkk(X)$ such that $\bigcap_{i=0}^n \Supp(D + \divisor(s_i)) = \Bs D = \emptyset$.
    Let $\phi_D: X \to \bbP_\kk^n = \bbP_\kk(H^0(X,D))$ be the morphism induced by $D$ (\cite[Chapter II, Theorem 7.1]{Har77}) and $Y$ be relative normalization (\Yang{}) of $X$ in $\bbP_\kk^n$.
    Then $\phi_D$ factors as $X \xrightarrow{f} Y \xrightarrow{g} \bbP_\kk^n$ with $g$ finite and $f$ a fibration.
    Let $A = g^* \calO_{\bbP_\kk^n}(1)$, then $A$ is ample since $g$ is finite, and $D \sim f^* A$ since $\phi_D^* \calO_{\bbP_\kk^n}(1) \sim D$.
    Hence we get the existence of the desired fibration.

    \begin{step}
      Uniqueness.
    \end{step}

    Let $f:X \to Y$ be a fibration and $A$ be an ample divisor on $Y$ such that $f^*D \sim_\bbQ A$.
    Replace $D,A$ by sufficiently divisible multiples, we may assume that $A$ is a very ample Cartier divisor and $f^*A \sim D$.
    Choose $s_1,\cdots, s_n \in H^0(Y,A)$ determining a closed immersion $\psi_n: Y \to \bbP^{n-1}$.
    For any $u_{n+1},\cdots,u_{n+m} \in H^0(X,D)$, let $\phi_{n+m}: X \to \bbP^{n+m-1}$ be the morphism associated to $f^*s_1,\cdots,f^*s_{n},u_{n+1},\cdots,u_{n+m}$.
    Let $L_n \subset \bbP^{n+m-1}$ be the linear subspace defined by $x_0=x_1=\cdots =x_n = 0$.
    Then $\phi_{n+m}(X) \cap L_n = \emptyset$ since $f^*s_0,\cdots,f^*s_{n}$ globally generate $\calO_X(D)$.
    We have a commutative diagram
    \[ \begin{tikzcd}
      X \ar[r,"\phi_{n+m}"] \ar[dr,"f"'] &\phi_{n+m}(X) \ar[d] \ar[r, hook] & \bbP^{n+m-1} \setminus L_n \ar[d, "\pi_n"] \\
        & Y \ar[r, "\psi_n", hook] & \bbP^{n-1}
    \end{tikzcd} \]
    where $\pi_n$ is the projection by $[x_1:\cdots:x_{n+m}] \mapsto [x_1:\cdots:x_n]$.
    Note that $\bbP^{n+m} \setminus L_n \cong \bbP^{n-1} \times \bbA^m$ and $\pi$ is the projection to $\bbP^{n-1}$ under this identification.
    Then $\pi_n$ and hence $\pi_n|_{\phi_{n+m}(X)}:\phi_{n+m}(X) \to Y$ is affine.
    Note that $X$ is proper.
    Hence $\pi_n|_{\phi_{n+m}(X)}$ is finite since it is affine and proper.
    Then it must be an isomorphism since $f$ is a fibration.

    Let $f':X \to Y'$ and $A'$ on $Y'$ be another fibration and ample divisor such that $f'^*A' \sim_\bbQ D$.
    Replace $D,A,A'$ by a sufficiently divisible multiple, we may assume that both $A$ and $A'$ are very ample and $f'^*A' \sim D$.
    Choose $t_1,\cdots, t_m \in H^0(Y',A')$ determining a closed immersion $\psi_m: Y' \to \bbP^{m-1}$ and let $u_i = f'^*t_i$ in above construction.
    The we have the commutative diagram 
    \[ \begin{tikzcd}
        & Y' \ar[r, "\psi_m", hook] & \bbP^{m-1}\\
      X \ar[r,"\phi_{n+m}"] \ar[dr,"f"'] \ar[ur,"f'"] &\phi_{n+m}(X) \ar[d] \ar[u] \ar[r, hook] & \bbP^{n+m-1} \setminus (L_n\cup L_m) \ar[d, "\pi_n"] \ar[u, "\pi_m"] \\
        & Y \ar[r, "\psi_n", hook] & \bbP^{n-1},
    \end{tikzcd} \]
    where $\pi_m$ and $L_m$ is defined similarly to $\pi_n$ and $L_n$.
    The middle column gives us desired isomorphism $h:Y \to Y'$.
    The left thing is to check $A \sim_\bbQ h^*A'$.
    Since $f^*A \sim f^*h^*A'=f'^*A'$, we have $f_*f^*\calO(A) \cong f_*f^*\calO(h^*A')$ and hence $\calO(A) \cong \calO(h^*A')$ by projection formula and the fact $f$ is a proper fibration.
  \end{proof}
  \begin{remark}
    The key point in the proof is to deduce the finiteness of $\pi|_{\phi_{n+m}(X)}$ by its affinity and properness.
    Here we essentially use the fact $X$ is proper.
    If $X$ is not proper, then the fibration associated to a semiample divisor is not necessarily unique.
    For example, let $D$ be any divisor on $\bbA^n_\kk$.
    Since $\Pic(\bbA^n) = 0$, $D \sim 0$ is base point free.
    Then any surjective linear projection $\bbA^n \to \bbA^m$ satisfying the condition in \cref{thm:semi-ample_fibration}.
  \end{remark}

  We can detect the dimension of the image of the semiample fibration by some asymptotic behavior of the dimension of global sections of multiples of $D$.

  \begin{definition}
    Let $D$ be a divisor on a projective variety $X$.
    We define the \emph{Iitaka dimension} of $D$ as
    \[ \kappa(X, D) \coloneqq \max \left\{ k \in \bbR_{\geq 0} \;\middle|\; \limsup_{m \to +\infty} \frac{h^0(X, mD)}{m^k} > 0 \right\} \in \{-\infty\} \cup \bbR_{\geq 0}. \]
    If $D = K_X$ is the canonical divisor, we call $\kappa(X, K_X)$ the \emph{Kodaira dimension} of $X$ and denote it by $\kappa(X)$ simply.
  \end{definition}

  \begin{definition}
    Let $D$ be a divisor on a projective variety $X$.
    We say that $D$ is \emph{big} if $\kappa(X, D) = \dim X$.
  \end{definition} 

  \begin{proposition}
    Let $D$ be an effective divisor on a projective variety $X$ with Iitaka fibration $f: X \setminus \bfB(D) \to Y$.
    Then we have $\kappa(X, D) = \dim Y$.
  \end{proposition}
  \begin{proof}
    \Yang{}
  \end{proof}

  \begin{remark}
    Iitaka dimension is not a numerical invariant, i.e. it may happen that $D \equiv D'$ but $\kappa(X, D) \neq \kappa(X, D')$.
    \Yang{arXiv:1904.10832}
  \end{remark}




    \Yang{}



\subsection{Base loci of numerical divisor classes}

  Another important property of base loci is intersection.
  For this purpose, we need that $X$ to be projective and we can only consider the numerical equivalence class of $D$.

  % \Yang{} It was first introduced by 

  \begin{definition}
    Let $D$ be a $\bbR$-divisor on a projective variety $X$.
    The \emph{augmented base locus} (also called the \emph{non-ample locus}) of $D$ is defined as
    \[ \bfB_+(D) \coloneqq \bigcap_{\substack{A \in \Div(X)_\bbR \text{ ample } \\ D - A \in \Div(X)_\bbQ}} \bfB(D - A). \]
    The \emph{restricted base locus} (also called the \emph{non-nef locus}) of $D$ is defined as
    \[ \bfB_-(D) \coloneqq \bigcup_{\substack{A \in \Div(X)_\bbR \text{ ample } \\ D + A \in \Div(X)_\bbQ}} \bfB(D + A). \]
    % \Yang{}
  \end{definition}

  \begin{remark}
    There seems many different definitions of $\bfB_+(D)$ and $\bfB_-(D)$ in the literature, see \cite{Birkar-2017-augmented-base-locus,Ein-Lazarsfeld-Mustața-Nakamaye-Popa-2006-asymptotic-invariants-base-loci}.
    The equivalence of these definitions \Yang{right?}
  \end{remark}

  \begin{example}
    Let $X = \Bl_p \bbP^2$ be the blow-up of $\bbP^2$ at a point $p$, and let $E$ be the exceptional divisor.
    Then for $a,b \in \bbR_{\geq 0}$, we have
    \[ \bfB_+(a\tilde{H} + bE) = \begin{cases}
      \emptyset & a > b, \\
      E & a \leq b,
    \end{cases} \quad \bfB_-(a\tilde{H} + bE) = \begin{cases}
      \emptyset & a > b, \\
      E & a < b.
    \end{cases} \]
    In particular, $\bfB_+(bE) = X$ but $\bfB_-(bE) = E$ for any $b > 0$.
     \Yang{}
  \end{example}

  \begin{remark}
    The restricted base locus $\bfB_-(D)$ is a countable union of closed subsets of $X$, so it is not necessarily closed.
     \Yang{arXiv:1212.3738}
  \end{remark}

  \begin{question}
    Suppose that $X$ is defined over a finitely generated field over $\bbQ$.
    Suppose that $D$ is pseudo-effective.
    Is $\bfB_-(D)$ always adelic closed or contained in a proper adelic closed subset of $X$?
     \Yang{}
  \end{question}

  \begin{example}
    \Yang{}
  \end{example}

  \begin{theorem}\label{prop:properties_augmented_and_restricted_base_loci}
    Let $D$ be a $\bbR$-divisor on a projective variety $X$.
    \begin{enumerate}
      \item Both restricted and augmented base loci only depend on the numerical equivalence class of $D$.
      \item We have $\bfB_-(D) \subseteq \bfB(D) \subseteq \bfB_+(D)$ when $D$ is a $\bbQ$-divisor.
      \item The intersection in the definition of $\bfB_+(D)$ will stabilize, i.e. fix a norm on $\upN^1(X)$, there exists $\varepsilon > 0$ such that for any ample $\bbQ$-divisor $A$ with $0 < \|A\| \leq \varepsilon$, we have $\bfB_+(D) = \bfB(D - A)$.
      %  \Yang{}
      \item Let $C$ be a curve on $X$ such that $C \not\subset \bfB_-(D)$, then $D \cdot C \geq 0$.
        In particular, $D$ is nef if and only if $\bfB_-(D) = \emptyset$; $D$ is pseudo-effective if and only if $\bfB_-(D) \neq X$.
       \Yang{}
      \item Augmented base locus $\bfB_+(D)$ coincides with the its exceptional locus, i.e.
      \[ \bfB_+(D) = \bigcup_{\substack{V \subseteq X \text{ subvariety} \\ D|_V \text{ not big}}} V. \]
      In particular, $D$ is big if and only if $\bfB_+(D) \neq X$ and $D$ is ample if and only if $\bfB_+(D) = \emptyset$. ref. \cite{Birkar-2017-augmented-base-locus}.
      %  \Yang{}
    \end{enumerate}
    \Yang{}
  \end{theorem}
  \begin{proof}
    \Yang{}
  \end{proof}

  \Yang{}



